\begin{mistake}{2026-04-12}{08}

\begin{problemcontent}
设
\[
A=\begin{bmatrix}
1&2&1\\
2&3&a+2\\
1&a&-2
\end{bmatrix},
\qquad
B=\begin{bmatrix}
1&1&a\\
-1&a&1\\
1&-1&2
\end{bmatrix},
\]
求使 \(A\) 与 \(B\) 不等价的 \(a\) 的值。

\end{problemcontent}

\begin{solutioncontent}
同型矩阵等价当且仅当秩相等，因此本题要求
\[
r(A)\ne r(B).
\]

先计算两个行列式。

对 \(A\) 按第一行展开：
\[
\begin{aligned}
|A|
&=\begin{vmatrix}3&a+2\\ a&-2\end{vmatrix}
-2\begin{vmatrix}2&a+2\\ 1&-2\end{vmatrix}
+\begin{vmatrix}2&3\\ 1&a\end{vmatrix}\\
&=( -6-a^2-2a)-2(-6-a-2)+(2a-3)\\
&=-a^2+2a+3\\
&=-(a-3)(a+1).
\end{aligned}
\]

对 \(B\) 按第一行展开：
\[
\begin{aligned}
|B|
&=\begin{vmatrix}a&1\\ -1&2\end{vmatrix}
-\begin{vmatrix}-1&1\\ 1&2\end{vmatrix}
+a\begin{vmatrix}-1&a\\ 1&-1\end{vmatrix}\\
&=(2a+1)-(-3)+a(1-a)\\
&=-a^2+3a+4\\
&=-(a-4)(a+1).
\end{aligned}
\]

若 \(|A|\ne 0\)，则 \(r(A)=3\)；若 \(|B|\ne 0\)，则 \(r(B)=3\)。
故只需检查使行列式为零的临界值
\[
a=3,\ 4,\ -1.
\]

当 \(a=3\) 时，\(|A|=0,\ |B|\ne 0\)，故
\[
r(A)<3,
\qquad
r(B)=3,
\]
于是 \(r(A)\ne r(B)\)。

当 \(a=4\) 时，\(|A|\ne 0,\ |B|=0\)，同理有 \(r(A)\ne r(B)\)。

当 \(a=-1\) 时，两者行列式都为 0，需继续验秩。

代入 \(a=-1\) 得
\[
A=\begin{bmatrix}
1&2&1\\
2&3&1\\
1&-1&-2
\end{bmatrix}.
\]
其二阶子式
\[
\begin{vmatrix}
1&2\\
2&3
\end{vmatrix}=-1\ne 0,
\]
故 \(r(A)=2\)。

又
\[
B=\begin{bmatrix}
1&1&-1\\
-1&-1&1\\
1&-1&2
\end{bmatrix}.
\]
其中前两行互为相反数，而
\[
\begin{vmatrix}
1&-1\\
1&2
\end{vmatrix}=3\ne 0,
\]
故 \(r(B)=2\)。

于是 \(a=-1\) 时仍有 \(r(A)=r(B)\)，两矩阵等价。

综上，所求为
\[
a=3\quad \text{或}\quad a=4.
\]

\end{solutioncontent}

\mistakeknowledge{
\begin{itemize}
  \item \textbf{核心公式/定理}：同型矩阵等价 \(\iff\) 秩相等。对三阶矩阵，先用行列式判断是否满秩，再对退化参数继续查二阶子式。
  \item \textbf{易错点}：公共零点 \(a=-1\) 不能直接判“不等价”。当两个行列式同时为 0 时，必须继续比较具体秩值。
  \item \textbf{关联知识}：判秩常用顺序是“先看最高阶子式，再降一阶检查”。这比盲目做大量行变换更稳、更快。
\end{itemize}}

\end{mistake}
